\section{Learning on the learner's cues: supervised one-step vs deep RL}\label{sec:learners}\label{sec:attrib}
This section gives the main result of the paper. Section~\ref{sec:results} showed that the non-oracle RL policy (NR2) beats every calibrated table and both hand-designed rules on the same cues, robustly across training seeds. Does that gain need reinforcement learning, or would a simpler learner on the same gNB-observable features capture it? Pre-registered controls (Amendments~\lit{8}, \lit{9} and \lit{10a}, Appendix~\ref{app:programme}) answer this in three steps: a supervised one-step learner with simulator labels (R1), the same learner with labels a gNB observes (R1b), a replication of the latter over four independent calibration draws and at two neighbouring operating points (R1c), and an RL policy with the learner's full action set (R6, Amendment~\lit{11}). The answer is that a supervised one-step model of transmission success, trained on observed ACK/NACK, beats the RL policy at every step, also when it is restricted to the RL policy's action set, and also when the RL policy is given the learner's full action set.

All controls act on the same Fixed base as NR2, see only the \lit{19} observable features of the policy state (Table~\ref{tab:info}; no speed class and no burst state) and are built on calibration traces only. They are evaluated on the \lit{500} NR2 confirmation traces, paired with common random numbers (Table~\ref{tab:attrib}). Margins are those of NR2 relative to the control, so a negative margin means that the control delivered more packets on time. At evaluation the RL policy acts greedily (largest logit), as the one-step learners do (largest predicted value).

\emph{Status and sequence.} These controls were not planned with the main round. HR1 had shown one-step headroom (Section~\ref{sec:hr1}) and NR2 had passed its confirmatory test when R1 and R2 were frozen (Amendment~\lit{8}). R1b (Amendment~\lit{9}) was designed after R1 had been seen, and R1c (Amendment~\lit{10a}) after R1b; the decision to make this comparison the main result of the paper was taken after R1b. R6 (Amendment~\lit{11}) was designed after an internal review of the draft had named the action-set difference as the main confound, and frozen before it ran. Each stage was frozen and hashed before it ran (Supplement~S3 gives the freeze and run times). R1 and R1b should therefore be read as exploratory results that R1c then tested. At M1, R1c evaluates on the NR2 confirmation traces, which NR2, R1, R1b and R2 also used; in the cells it uses the G10 evaluation traces. The R1c comparator in G2 includes five seeds re-evaluated under Deviation~D1, which was recorded together with Amendment~\lit{9}, before Amendment~\lit{10a}; Amendment~\lit{10a} specified the corrected evaluation. All three stages are in the project-wide count, and the NR2 comparisons with the ACK-trained learners keep $p=$~\NBonfPProjGbtResAck{} and \NBonfPProjGbtFullAck{} after the project-wide Bonferroni factor (Section~\ref{sec:protocol}).

\emph{What differs between the learners.} The comparison is not matched in every respect. The RL policy chooses a residual step of at most one CQI around its base (base $\pm\lit{1}$), is trained on-policy for \NCfgEpDl{} episodes per seed with a prior towards the base rule (a logit bias and a deviation penalty, Appendix~\ref{app:hp}), and was not tuned beyond one hyper-parameter set. The one-step learners are fitted offline to the logged calibration traces; GBT-res and GBT-res-ACK have the same action set as the RL policy, while GBT-full and GBT-full-ACK may choose any offset from $-\lit{4}$ to $+\lit{2}$. The restricted comparison removes the action-set difference on the learner's side; R6 removes it on the RL side (Section~\ref{sec:r6}).

\begin{table*}[t]
\caption{Learners and rules on the learner's cues, on the \lit{500} NR2 confirmation traces. On-time: mean on-time ratio (NR2: the \NCfgSeeds{} seeds pooled). NR2 margin: NR2 relative to the control (negative: the control is better), with the \NLevTrace\,\% CI, the hierarchical seed$\times$trace \NLevHier\,\% CI and the NR2 seeds above the control; Holm over the controls of each pre-registered family. Jain index over per-user on-time ratios and the mean of the bottom decile of per-user on-time ratios, from the exact-path audit; a user resolves only about \NPSPktUser{} packets per trace, so the bottom decile is coarse and often zero. --: not computed. $^\dagger$RL-full row: margin of RL-full relative to NR2 (positive: RL-full is better) and RL-full seeds above NR2; descriptive. RL-full against GBT-full-ACK: Table~\ref{tab:r6}.}\label{tab:attrib}
\centering\footnotesize
\setlength{\tabcolsep}{3pt}
\begin{tabular}{@{}lccccccc@{}}
\toprule
Scheme (labels / training data) & On-time & NR2 margin (\%) & \NLevTrace\,\% CI & hier.\ \NLevHier\,\% CI & NR2 seeds above & Jain & Bottom decile\\
\midrule
NR2: residual A2C, non-oracle base & \NPSOntNRtwo & -- & -- & -- & -- & \NPSJainNRtwo & \NPSBotNRtwo\\
Combined table (Context-joint-500) & \NPSOntCtxJ & +\NRtwo{margin_ctxj} & [\NRtwo{ci_lo_ctxj}, \NRtwo{ci_hi_ctxj}] & [\NRtwo{hier_lo}, \NRtwo{hier_hi}] & \NRtwo{seeds_above} & \NPSJainCtxJ & \NPSBotCtxJ\\
Designed rule (as written) & \NPSOntRule & \NPSRuleMar & [\NPSRuleLo, \NPSRuleHi] & [\NPSRuleHierLo, \NPSRuleHierHi] & \NPSRuleSeeds & -- & --\\
Designed rule (calibrated thresholds) & \NPSOntRuleCal & \NPSRuleCalMar & [\NPSRuleCalLo, \NPSRuleCalHi] & [\NPSRuleCalHierLo, \NPSRuleCalHierHi] & \NPSRuleCalSeeds & -- & --\\
GBT-res (true-SINR labels, Fixed logging) & \NPSOntGbtRes & \NPSGbtResMar & [\NPSGbtResLo, \NPSGbtResHi] & [\NPSGbtResHierLo, \NPSGbtResHierHi] & \NPSGbtResSeeds & \NPSJainGbtRes & \NPSBotGbtRes\\
GBT-full (true-SINR labels, Fixed logging) & \NPSOntGbtFull & \NPSGbtFullMar & [\NPSGbtFullLo, \NPSGbtFullHi] & [\NPSGbtFullHierLo, \NPSGbtFullHierHi] & \NPSGbtFullSeeds & \NPSJainGbtFull & \NPSBotGbtFull\\
\textbf{GBT-res-ACK} (observed ACK/NACK, exploration) & \NPSOntGbtResAck & \NPSGbtResAckMar & [\NPSGbtResAckLo, \NPSGbtResAckHi] & [\NPSGbtResAckHierLo, \NPSGbtResAckHierHi] & \NPSGbtResAckSeeds & \NPSJainGbtResAck & \NPSBotGbtResAck\\
\textbf{GBT-full-ACK} (observed ACK/NACK, exploration) & \NPSOntGbtFullAck & \NPSGbtFullAckMar & [\NPSGbtFullAckLo, \NPSGbtFullAckHi] & [\NPSGbtFullAckHierLo, \NPSGbtFullAckHierHi] & \NPSGbtFullAckSeeds & \NPSJainGbtFullAck & \NPSBotGbtFullAck\\
\midrule
RL-full: residual A2C, offsets $-\lit{4}\dots+\lit{2}$ (R6)$^\dagger$ & \NRsixOnt & \NRsixVsNRtwo & [\NRsixVsNRtwoLo, \NRsixVsNRtwoHi] & [\NRsixVsNRtwoHierLo, \NRsixVsNRtwoHierHi] & \NRsixVsNRtwoSeeds & -- & --\\
\bottomrule
\end{tabular}
\end{table*}

\subsection{A one-step learner with simulator labels (R1)}
The headroom analysis (Section~\ref{sec:hr1}) found that a gradient-boosted predictor of the expected useful bits of a first transmission leaves room over every calibrated table. The first control deploys such a predictor as a policy (R1). One model per candidate offset is trained on the calibration traces only, with the Fixed rule generating the logged states. The labels use the simulator's true SINR; they are training targets only, and at deployment the model sees the \lit{19} observable features. At each decision it chooses the offset with the largest predicted expected useful bits, within base $\pm\lit{1}$ (GBT-res, the action set of NR2) or over all offsets (GBT-full). Both beat NR2: NR2 delivers \NPSGbtResAdv\,\% fewer packets on time than GBT-res and \NPSGbtFullAdv\,\% fewer than GBT-full, and no NR2 seed is above either. The same holds on the M1 evaluation traces (NR2 margins \NPSGbtResMoneMar\,\% and \NPSGbtFullMoneMar\,\%).

\subsection{A one-step learner on observed ACK/NACK (R1b)}\label{sec:r1b}
A gNB never sees the true SINR, so the second control (R1b) uses only what it can observe: the ACK/NACK outcome of each first transmission. Outcomes of offsets other than the base are observed only if they are tried, so the calibration traces were logged with $\varepsilon$-greedy exploration ($\varepsilon=$~\NPSExplEps: half of the decisions use the base offset, the rest an offset drawn uniformly from $-\lit{4}\ldots+\lit{2}$). A gradient-boosted classifier with the same settings estimates $\hat p(\text{ACK}\mid x,\text{MCS})$ from the observable features and the MCS, and the policy chooses the offset with the largest $\hat p\times$ useful bits, again within base $\pm\lit{1}$ (GBT-res-ACK) or over all offsets (GBT-full-ACK). These versions, trained only on labels a gNB observes, beat NR2 by more: NR2 delivers \NPSGbtResAckAdv\,\% and \NPSGbtFullAckAdv\,\% fewer packets on time (hierarchical \NLevHier\,\% CIs of the NR2 margin [\NPSGbtResAckHierLo, \NPSGbtResAckHierHi]\,\% and [\NPSGbtFullAckHierLo, \NPSGbtFullAckHierHi]\,\%; \NPSGbtFullAckSeeds{} of \NCfgSeeds{} seeds above). Against the combined table they gain \NPSGbtResAckOverCtxJ\,\% and \NPSGbtFullAckOverCtxJ\,\%.

The success model is well calibrated: the expected calibration error is at most \NPSEceMax{} over the cross-fitted calibration predictions, the explored decisions only, and the decisions taken by both policies on the confirmation traces (Brier score \NPSBrierXfit{} cross-fitted). As an observation, not a pre-registered test, the ACK-trained policies also did better than their true-SINR twins (\NPSTwinGbtResAck\,\% and \NPSTwinGbtFullAck\,\%). The two differ in their labels \emph{and} in their logging policy (exploration against the Fixed rule), so this comparison does not isolate the effect of the labels.

\subsection{Replication across calibration draws and operating points (R1c)}\label{sec:r1c}
Each R1 and R1b policy comes from a single calibration draw at the main operating point. The pre-registered replication R1c (Amendment~\lit{10a}) refits both ACK-trained learners, with the same settings and the same exploration logging, on four independent calibration draws at M1, evaluated on the NR2 traces against NR2. In each neighbouring cell (G2, G3) it trains them on four fresh draws of the cell's own calibration traces and compares them with the cell's ten-seed non-oracle RL policy and with the cell's nine calibrated comparators. A trace set \emph{holds} if every fit beats or matches the RL policy (match: Holm not significant and $|\text{margin}|<\lit{0.5}$\,\%) and, in the cells, also beats all nine comparators. Table~\ref{tab:r1c} gives the result; per-fit margins, intervals and Holm $p$-values are in Supplement~S2.

\wordingsel{R1Cres}{R1C}{verdict_res}{unrun} \wordingsel{R1Cfull}{R1C}{verdict_full}{unrun} \ifitemrun{R1C}{No fit merely matched (match is a point-estimate rule, not an equivalence test): even the fit with the lowest on-time ratio beat the RL policy in every trace set. On the NR2 traces the RL margin over the weakest fit is $\Ronec{worstmargin_nrtwo_res}$\,\% for GBT-res-ACK and $\Ronec{worstmargin_nrtwo_full}$\,\% for GBT-full-ACK; fit~\lit{1} at M1 is the draw reported for R1b, so the M1 replication consists of that draw and three new ones; for GBT-full-ACK the R1b draw is the weakest of the four (fit~\Ronec{worstfit_nrtwo_full}), and all three new draws did better. In G2 the cell-trained GBT-full-ACK reaches \Ronec{mean_gtwo_full} on time on average against \Ronec{rl_gtwo} for the RL policy (weakest fit: $\Ronec{worstmargin_gtwo_full}$\,\%), and in G3, close to the ceiling, \Ronec{mean_gthree_full} against \Ronec{rl_gthree} ($\Ronec{worstmargin_gthree_full}$\,\%). The fits trained at M1 and applied unchanged in the cells did about as well (secondary; Table~\ref{tab:r1c}).}{}

\begin{table*}[t]
\caption{Replication of the ACK-trained one-step learner (R1c, Amendment~\lit{10a}). Four independent calibration draws per trace set; on-time: mean over the four fits [lowest, highest fit]; state: holds (all four fits beat or match the RL policy and, in the cells, beat all nine calibrated comparators), fails (none does) or split; worst fit: the fit with the lowest on-time ratio, the RL margin over it (\%; negative: the learner is better) and its outcome (LOSES: the RL policy loses to the learner). RL: NR2 on the NR2 traces, the NR-G10 policy (\NCfgSeeds{} seeds; G2 seeds \lit{6}--\lit{10} corrected, Deviation~D1) in the cells. Transfer: the M1 fits applied unchanged in the cells (secondary; the state uses the same rule but carries no pre-registered claim). Per-fit intervals and Holm $p$-values: Supplement~S2.}\label{tab:r1c}
\centering\footnotesize
\setlength{\tabcolsep}{2.5pt}
\resizebox{\textwidth}{!}{%
\begin{tabular}{@{}llcccccc@{}}
\toprule
Trace set & RL on-time & \multicolumn{3}{c}{GBT-res-ACK} & \multicolumn{3}{c}{GBT-full-ACK}\\
\cmidrule(lr){3-5}\cmidrule(lr){6-8}
 & & state & on-time [min, max] & worst fit & state & on-time [min, max] & worst fit\\
\midrule
NR2 traces, M1 fits & \Ronec{rl_nrtwo} & \Ronec{state_nrtwo_res} & \Ronec{mean_nrtwo_res} [\Ronec{min_nrtwo_res}, \Ronec{max_nrtwo_res}] & \Ronec{worstfit_nrtwo_res}: \Ronec{worstmargin_nrtwo_res}, \Ronec{worstout_nrtwo_res} & \Ronec{state_nrtwo_full} & \Ronec{mean_nrtwo_full} [\Ronec{min_nrtwo_full}, \Ronec{max_nrtwo_full}] & \Ronec{worstfit_nrtwo_full}: \Ronec{worstmargin_nrtwo_full}, \Ronec{worstout_nrtwo_full}\\
G2, cell fits & \Ronec{rl_gtwo} & \Ronec{state_gtwo_res} & \Ronec{mean_gtwo_res} [\Ronec{min_gtwo_res}, \Ronec{max_gtwo_res}] & \Ronec{worstfit_gtwo_res}: \Ronec{worstmargin_gtwo_res}, \Ronec{worstout_gtwo_res} & \Ronec{state_gtwo_full} & \Ronec{mean_gtwo_full} [\Ronec{min_gtwo_full}, \Ronec{max_gtwo_full}] & \Ronec{worstfit_gtwo_full}: \Ronec{worstmargin_gtwo_full}, \Ronec{worstout_gtwo_full}\\
G3, cell fits & \Ronec{rl_gthree} & \Ronec{state_gthree_res} & \Ronec{mean_gthree_res} [\Ronec{min_gthree_res}, \Ronec{max_gthree_res}] & \Ronec{worstfit_gthree_res}: \Ronec{worstmargin_gthree_res}, \Ronec{worstout_gthree_res} & \Ronec{state_gthree_full} & \Ronec{mean_gthree_full} [\Ronec{min_gthree_full}, \Ronec{max_gthree_full}] & \Ronec{worstfit_gthree_full}: \Ronec{worstmargin_gthree_full}, \Ronec{worstout_gthree_full}\\
G2, transfer & \Ronec{rl_gtwo} & \Ronec{state_gtwotr_res} & \Ronec{mean_gtwotr_res} [\Ronec{min_gtwotr_res}, \Ronec{max_gtwotr_res}] & \Ronec{worstfit_gtwotr_res}: \Ronec{worstmargin_gtwotr_res}, \Ronec{worstout_gtwotr_res} & \Ronec{state_gtwotr_full} & \Ronec{mean_gtwotr_full} [\Ronec{min_gtwotr_full}, \Ronec{max_gtwotr_full}] & \Ronec{worstfit_gtwotr_full}: \Ronec{worstmargin_gtwotr_full}, \Ronec{worstout_gtwotr_full}\\
G3, transfer & \Ronec{rl_gthree} & \Ronec{state_gthreetr_res} & \Ronec{mean_gthreetr_res} [\Ronec{min_gthreetr_res}, \Ronec{max_gthreetr_res}] & \Ronec{worstfit_gthreetr_res}: \Ronec{worstmargin_gthreetr_res}, \Ronec{worstout_gthreetr_res} & \Ronec{state_gthreetr_full} & \Ronec{mean_gthreetr_full} [\Ronec{min_gthreetr_full}, \Ronec{max_gthreetr_full}] & \Ronec{worstfit_gthreetr_full}: \Ronec{worstmargin_gthreetr_full}, \Ronec{worstout_gthreetr_full}\\
\bottomrule
\end{tabular}}
\end{table*}

\subsection{RL with the full action set (R6)}\label{sec:r6}
GBT-full-ACK may choose offsets that NR2 cannot reach. R6 (Amendment~\lit{11}) trained the NR2 recipe with one change, a residual step $\delta\in\{-\lit{2},\dots,+\lit{4}\}$ on the Fixed $-\lit{2}$ base, so that the RL policy (RL-full) chooses from the same offsets $-\lit{4}\dots+\lit{2}$ as GBT-full-ACK. The prior towards the base rule was kept, ten fresh seeds were trained for \NCfgEpDl{} episodes each, and the policy was evaluated greedily on the NR2 traces against each of the four GBT-full-ACK draws (Holm over four). The pre-registered verdict is \emph{the supervised learner beats RL at both action sets}: With the same action set as the one-step learner (offsets $-\lit{4}\dots+\lit{2}$), the RL policy still delivered fewer packets on time than the one-step learner trained on observed ACK/NACK, for each of the four calibration draws: the supervised learner beats RL at both action sets.

RL-full reaches an on-time ratio of \NRsixOnt. It delivers \NRsixFullLossMin--\NRsixFullLossMax\,\% fewer packets on time than the four draws (Table~\ref{tab:r6}); every hierarchical interval lies below zero, no seed is ahead of any draw, and the comparison keeps $p=$~\NBonfPProjRsix{} after the project-wide Bonferroni factor. It also loses to all four draws of GBT-res-ACK, by \NRsixResLossMin--\NRsixResLossMax\,\% (secondary). The wider set was used: on average \NRsixOutside\,\% of its decisions took a step outside $\pm\lit{1}$, most of them the largest step ($\delta=+\lit{4}$, \NRsixPlusFour\,\%). It gained little from them. Against NR2 its margin is \NRsixVsNRtwo\,\% (\NLevTrace\,\% CI [\NRsixVsNRtwoLo, \NRsixVsNRtwoHi]\,\%), which is not robust across seeds (hierarchical \NLevHier\,\% CI [\NRsixVsNRtwoHierLo, \NRsixVsNRtwoHierHi]\,\%, \NRsixVsNRtwoSeeds{} of \NCfgSeeds{} seeds above; descriptive), so it closes \NRsixClosure\,\% of the gap between NR2 and GBT-full-ACK. Against the combined table it keeps the RL advantage (\NRsixVsCtxJ\,\%, hierarchical \NLevHier\,\% CI [\NRsixVsCtxJHierLo, \NRsixVsCtxJHierHi]\,\%, \NRsixVsCtxJSeeds{} of \NCfgSeeds{} seeds above). The action-set difference therefore does not explain why the one-step learner wins.

\begin{table}[t]
\caption{RL with the full action set (R6, Amendment~\lit{11}) against the four draws of GBT-full-ACK on the NR2 traces. Margin: RL-full relative to the draw (negative: the one-step learner is better), with the hierarchical seed$\times$trace \NLevHier\,\% CI; Holm over the four draws. RL-full seeds above any draw: \NRsixFullSeeds{} of \NCfgSeeds. Draw~\lit{1} is the R1b model.}\label{tab:r6}
\centering\footnotesize
\begin{tabular}{@{}lcc@{}}
\toprule
GBT-full-ACK draw & Margin (\%) & hier.\ \NLevHier\,\% CI\\
\midrule
\lit{1} & \NRsixFullOne & [\NRsixFullOneHierLo, \NRsixFullOneHierHi]\\
\lit{2} & \NRsixFullTwo & [\NRsixFullTwoHierLo, \NRsixFullTwoHierHi]\\
\lit{3} & \NRsixFullThree & [\NRsixFullThreeHierLo, \NRsixFullThreeHierHi]\\
\lit{4} & \NRsixFullFour & [\NRsixFullFourHierLo, \NRsixFullFourHierHi]\\
\bottomrule
\end{tabular}
\end{table}

An exploratory diagnostic on the role of label type (R6b) was inconclusive; the diagnosis (\path{Optimized/v17x/Dataset/presub/r6/R6b_diagnosis.md}) ruled out pipeline errors but did not explain why a model trained on exact success probabilities from the same exploration logs underperformed the ACK-trained model. It is not reported and is left for future work.

\subsection{Side effects}
The ACK-trained learners run at a first-transmission block-error rate of \NPSBlerGbtResAck\,\% (GBT-res-ACK) and \NPSBlerGbtFullAck\,\% (GBT-full-ACK), against \NPSBlerNRtwo\,\% for NR2 and \NPSBlerCtxJ\,\% for the combined table; the true-SINR learners run at \NPSBlerGbtRes\,\% and \NPSBlerGbtFull\,\%. All of them sit well above the OLLA target, as is consistent with policies that accept more block errors for larger transport blocks under deadline pressure. Their Jain index over per-user on-time ratios is at least that of NR2 (Table~\ref{tab:attrib}). We make no claim about tail users: the pre-registered edge5 guard was not computed for these controls, and with about \NPSPktUser{} packets per user and trace the bottom decile of per-user on-time ratios is near zero for every scheme and cannot separate them.

\subsection{What the one-step learner does not do}\label{sec:myopic}
The one-step learner is myopic. It scores each first transmission by its predicted success probability times its useful bits (the payload bits a successful transmission delivers) and gives no credit for what a decision does to later HARQ rounds, queues or deadlines. The RL policy is trained on exactly that sequential credit, and the myopic model still wins in every trace set. This should be read with the structure of the tested points in mind. With a budget of \NCfgDeadline{} TTIs and an HARQ round-trip time of \NCfgRtt{} TTIs, a packet can receive at most one retransmission before its deadline, so the first transmission decides most outcomes and the problem is close to a contextual bandit, which favours a one-step learner. The structure does not explain the whole result, however: in G3 the deadline is \NGthreeDeadline{} TTIs, so two retransmissions fit, and the ACK-trained learners still beat the RL policy there (\Ronec{mean_gthree_res} and \Ronec{mean_gthree_full} on time against \Ronec{rl_gthree}). The result is consistent with the headroom analysis (Section~\ref{sec:hr1}): most of what there is to learn here appears to be a per-decision trade-off between success probability and transport-block size, which a supervised model estimates directly from labelled outcomes while the RL policy must find it through a noisy, delayed return. It also shows that HR1 is not an upper bound on what a one-step learner realises over a trace: the ACK-trained learners gain \NPSGbtResAckOverCtxJ\,\% and \NPSGbtFullAckOverCtxJ\,\% over the combined table, at least as much as the one-step estimates in different units. We have not tested this reading. R6 has removed one alternative explanation, the action set (Section~\ref{sec:r6}). Two runs would separate the remaining ones: an RL policy (or fitted $Q$-iteration) trained on the same exploration logs with and without discounting, and a residual RL policy on top of the one-step learner. A point where several HARQ rounds fit before the deadline would test whether sequential credit then pays off.
